\chapter{Studi Literatur}
\label{ch:studi-literatur}

\section{Rekam Medis Elektronik}
\label{sec:rme}

\subsection{Definisi dan Regulasi RME di Indonesia}
\label{subsec:definisi-regulasi-rme}

Sesuai dengan Pasal 1 ayat (1) dan (2) Peraturan Menteri Kesehatan Republik Indonesia Nomor 24 Tahun 2022 tentang Rekam Medis, rekam medis adalah dokumen yang berisikan data identitas pasien, pemeriksaan, pengobatan, tindakan, dan pelayanan lain yang telah diberikan kepada pasien, sedangkan Rekam Medis Elektronik (RME) adalah rekam medis yang dibuat dengan menggunakan sistem elektronik yang diperuntukkan bagi penyelenggaraan rekam medis (\cite{permenkes24_2022}). Berdasarkan definisi tersebut, RME bukan sekadar bentuk digital dari berkas rekam medis konvensional, melainkan mencakup keseluruhan proses penyelenggaraan yang berkelanjutan, mulai dari bagaimana data dibuat, diakses, disimpan, hingga didistribusikan kepada pihak-pihak tertentu untuk memenuhi kebutuhan pelayanan kesehatan. Permenkes Nomor 24 Tahun 2022 juga mewajibkan seluruh fasilitas pelayanan kesehatan (fasyankes) di Indonesia, mulai dari praktik mandiri tenaga kesehatan hingga rumah sakit, untuk menyelenggarakan RME sebagai pengganti rekam medis konvensional (\cite{permenkes24_2022}).

Selain mewajibkan bentuk elektronik, Permenkes Nomor 24 Tahun 2022 turut mengatur hak pasien atas data rekam medis miliknya, termasuk larangan akses oleh pihak yang tidak berwenang (\cite{permenkes24_2022}). Ketentuan ini menjadi dasar regulasi bahwa setiap akses terhadap data RME seorang pasien pada dasarnya harus melalui mekanisme yang dapat mempertanggungjawabkan siapa yang mengakses, kapan, dan untuk keperluan apa. Hal tersebut sejalan dengan Undang-Undang Nomor 19 Tahun 2016 tentang Perubahan atas Undang-Undang Nomor 11 Tahun 2008 tentang Informasi dan Transaksi Elektronik, yang mengategorikan RME sebagai dokumen elektronik sehingga harus dibuat dan dikelola dalam sistem elektronik dengan jaminan keamanan dan akuntabilitas agar dapat digunakan sebagai alat bukti yang sah secara hukum (\cite{uu19_2016}). Dengan kata lain, regulasi di Indonesia tidak hanya menuntut RME diselenggarakan secara elektronik, tetapi juga menuntut adanya mekanisme yang menjamin kerahasiaan dan keutuhan data tersebut sepanjang siklus penyelenggaraannya, sebuah kebutuhan yang menjadi dasar bagi berbagai penelitian pengembangan sistem manajemen RME yang berfokus pada keamanan data, salah satunya DecMed (\cite{DecMed}).

\subsection{Kebutuhan Keamanan dan Akuntabilitas Data RME}
\label{subsec:keamanan-akuntabilitas-rme}

Data RME tergolong sebagai data yang sensitif karena memuat informasi identitas dan kondisi kesehatan pasien, sehingga keamanan dan kerahasiaannya menjadi perhatian utama dalam berbagai penelitian sistem informasi kesehatan. Tinjauan sistematis terhadap implementasi RME menunjukkan bahwa aspek keamanan dan kerahasiaan data medis pasien secara konsisten menjadi isu sentral, dan kegagalan menjaga aspek tersebut berdampak langsung pada menurunnya kepercayaan pasien terhadap sistem pelayanan kesehatan yang digunakan (\cite{pramesti2024keamanan}). Kerahasiaan data RME tidak dapat dijaga hanya dengan membatasi siapa saja yang boleh mengakses data (kontrol akses), tetapi juga membutuhkan mekanisme yang dapat menunjukkan bahwa pembatasan tersebut benar-benar dipatuhi dan dapat ditelusuri kembali ketika terjadi insiden.

Kebutuhan penelusuran kembali (\textit{traceability}) inilah yang menjadikan akuntabilitas sebagai aspek keamanan yang tidak dapat dipisahkan dari kerahasiaan dan integritas data RME. Akuntabilitas dalam konteks ini merujuk pada kemampuan suatu sistem untuk menghubungkan setiap tindakan yang dilakukan terhadap data dengan identitas pihak yang melakukannya, sehingga setiap pihak yang mengakses maupun mengubah data RME dapat dimintai pertanggungjawaban atas tindakannya. Sebagai contoh, kerangka kerja DecMed yang menjadi basis pengembangan pada Tugas Akhir ini dirancang agar seluruh akses terhadap data RME hanya dapat dilakukan atas persetujuan aktif pasien melalui mekanisme kontrol akses berbasis \textit{capability} (\cite{DecMed}). Namun demikian, tanpa mekanisme yang menjamin akuntabilitas, kontrol akses yang ketat sekalipun tidak dapat membuktikan bahwa suatu pelanggaran keamanan benar-benar terjadi, siapa pihak yang bertanggung jawab, maupun bagaimana kronologi pelanggaran tersebut (\cite{DecMed}). Kebutuhan inilah yang menjadi latar belakang pentingnya mekanisme audit trail sebagai pelengkap dari kontrol akses pada sistem manajemen RME, yang akan dibahas lebih lanjut pada Subbab~\ref{sec:audit-trail}.


% Input
\section{Audit Trail}
\label{sec:audit-trail}

\subsection{Definisi dan Konsep Audit Trail}
\label{subsec:definisi-audit-trail}

Audit trail didefinisikan sebagai rekaman kronologis atas aktivitas sistem yang cukup untuk memungkinkan rekonstruksi, peninjauan, dan pemeriksaan urutan kejadian yang mengelilingi atau menyebabkan suatu operasi, prosedur, atau peristiwa dalam suatu transaksi, sejak awal hingga hasil akhirnya (\cite{iso27789_2021}). Definisi ini menegaskan bahwa audit trail bukan sekadar kumpulan catatan yang berdiri sendiri, melainkan suatu rangkaian catatan yang secara kolektif membentuk kronologi yang dapat ditelusuri ulang, sehingga setiap catatan di dalamnya perlu memiliki keterkaitan yang jelas dengan catatan lain dari kejadian yang sama. Standar yang sama juga membedakan secara eksplisit antara \textit{audit record} sebagai catatan satu peristiwa tunggal dalam siklus hidup suatu rekam medis elektronik, \textit{audit log} sebagai kumpulan audit record yang tersusun secara kronologis, dan \textit{audit trail} sebagai kumpulan audit log yang secara keseluruhan membentuk riwayat lengkap suatu entitas (\cite{iso27789_2021}).

Konsep audit trail perlu dibedakan secara tegas dari konsep kontrol akses maupun pencatatan izin akses (\textit{access grant log}). Audit trail bersifat komplementer terhadap kontrol akses, bukan pengganti maupun bagian dari kontrol akses itu sendiri: kontrol akses menentukan kebijakan mengenai siapa \textit{berhak} mengakses data, sedangkan audit trail mencatat apa yang \textit{benar-benar terjadi} pada sistem, termasuk apabila terjadi penyimpangan dari kebijakan tersebut (\cite{iso27789_2021}). Perbedaan ini penting untuk digarisbawahi karena suatu sistem dapat saja memiliki kontrol akses yang ketat dan mencatat metadata pemberian izin secara lengkap, namun tetap tidak memiliki audit trail yang memadai apabila tidak ada catatan terpisah mengenai peristiwa akses aktual (\textit{access event}) yang terjadi setelah izin tersebut diberikan. Dengan kata lain, audit trail idealnya mencatat peristiwa (\textit{event}) sebagai representasi dari tindakan nyata yang terjadi pada sistem, bukan semata representasi dari kebijakan atau otorisasi yang berlaku.

\subsection{Karakteristik Audit Trail yang Baik}
\label{subsec:karakteristik-audit-trail}

Agar dapat berfungsi sebagai bukti yang dapat diandalkan, audit trail perlu memenuhi sejumlah karakteristik tertentu. Karakteristik pertama adalah kemampuan mengidentifikasi pengguna secara tidak ambigu (\textit{unambiguous identification}), yaitu audit trail harus menyediakan data yang cukup untuk mengidentifikasi secara pasti setiap pengguna maupun sistem eksternal yang telah mengakses atau menerima data rekam medis (\cite{iso27789_2021}). Karakteristik kedua adalah integritas (\textit{integrity}) dan kerahasiaan (\textit{confidentiality}) dari audit trail itu sendiri, yaitu jaminan bahwa catatan yang telah dibuat terlindungi dari upaya perusakan (\textit{tampering}), termasuk perlunya audit trail mencatat setiap tindakan yang dilakukan terhadap audit trail itu sendiri, seperti kapan sistem audit tidak berfungsi atau dinonaktifkan (\cite{iso27789_2021}).

Karakteristik lain yang tidak kalah penting adalah kelengkapan (\textit{completeness}) dan ketersediaan (\textit{availability}). Kelengkapan berarti audit trail harus mencatat seluruh tindakan yang relevan terhadap data -- termasuk pembuatan, pembacaan, pembaruan, maupun pengarsipan -- setiap kali tindakan tersebut dilakukan terhadap data kesehatan pribadi (\cite{iso27789_2021}), sedangkan ketersediaan berarti sistem audit harus memastikan entri tercatat setiap kali sistem informasi kesehatan sedang beroperasi (\cite{iso27789_2021}). Selain keempat karakteristik tersebut, audit trail yang baik juga perlu mendukung kemampuan deteksi terhadap potensi penyalahgunaan oleh pihak internal (\textit{insider threat}), mengingat pelaku insider pada umumnya memiliki akses sah terhadap sistem sehingga satu-satunya cara untuk mendeteksi maupun membuktikan penyalahgunaan tersebut adalah melalui pencatatan aktivitas yang menyeluruh dan tidak dapat dimanipulasi oleh pelaku itu sendiri (\cite{cert2022insiderthreat}). Karakteristik-karakteristik tersebut menjadi acuan utama dalam menilai apakah suatu mekanisme pencatatan pada sistem informasi, termasuk pada DecMed, telah dapat disebut sebagai audit trail yang memadai.

\subsection{Komponen Audit Trail}
\label{subsec:komponen-audit-trail}

Untuk dapat berfungsi secara menyeluruh, suatu sistem audit trail perlu dibangun melalui beberapa komponen fungsional yang saling berurutan. Kerangka kerja audit trail untuk rekam medis elektronik membagi persoalan menjadi penentuan peristiwa apa saja yang perlu dipicu untuk dicatat (\textit{trigger events}), format dan isi dari setiap catatan yang dihasilkan, hingga bagaimana catatan tersebut dikelola secara aman sepanjang siklus hidupnya (\cite{iso27789_2021}). Pembagian yang serupa juga ditemukan pada kerangka manajemen log komputer secara umum, yang memandang manajemen log sebagai suatu infrastruktur yang terdiri atas perangkat keras, perangkat lunak, jaringan, dan media yang digunakan untuk menghasilkan (\textit{generate}), mengirimkan (\textit{transmit}), menyimpan (\textit{store}), serta menganalisis (\textit{analyze}) data log (\cite{nist80092}). Berdasarkan kedua kerangka tersebut, komponen audit trail pada Tugas Akhir ini dikelompokkan menjadi empat bagian, yaitu penentuan event, pengumpulan event, penyimpanan event, dan pembacaan event, sebagaimana diuraikan pada subbagian berikut.

\subsubsection{Penentuan Event}
\label{subsubsec:penentuan-event}

Komponen pertama adalah penentuan peristiwa (\textit{event}) apa saja yang perlu memicu pencatatan audit trail. Peristiwa pemicu (\textit{trigger events}) pada sistem rekam medis elektronik pada umumnya ditentukan berdasarkan skala, tujuan, serta kebijakan privasi dan keamanan dari masing-masing sistem informasi kesehatan, dengan cakupan minimal berupa peristiwa akses (\textit{access events}) dan peristiwa kueri (\textit{query events}) terhadap data kesehatan pribadi (\cite{iso27789_2021}). Standar tersebut turut memberikan contoh peristiwa yang secara eksplisit berada di luar cakupan audit trail rekam medis, seperti peristiwa mulai/berhenti aplikasi, peristiwa autentikasi pengguna, maupun peristiwa yang dihasilkan oleh sistem operasi (\cite{iso27789_2021}).

Di luar peristiwa akses dan kueri, penentuan trigger events pada suatu sistem juga dapat diturunkan melalui pendekatan pemodelan ancaman (\textit{threat modeling}), seperti metode STRIDE, yang mengidentifikasi ancaman keamanan pada suatu sistem secara sistematis dan dapat digunakan sebagai dasar penentuan kebutuhan pencatatan (\cite{laksono2021stride}).

\subsubsection{Konten Wajib dan Konten Tambahan pada Event}
\label{subsubsec:konten-event}

Setiap audit record, terlepas dari jenis peristiwa apa yang dicatatnya, memuat sekumpulan bidang data minimal yang bersifat wajib (\textit{mandatory}). Terdapat sembilan bidang yang wajib diisi pada setiap event apa pun jenisnya, yaitu identitas unik peristiwa (\textit{EventID}), jenis tindakan yang dilakukan (\textit{EventActionCode}), waktu terjadinya peristiwa (\textit{EventDateTime}), identitas pengguna atau proses pelaku tindakan (\textit{UserID}), identitas sumber audit atau sistem yang menghasilkan catatan tersebut (\textit{AuditSourceID}), serta empat bidang yang secara bersama-sama mengidentifikasi objek data yang menjadi sasaran tindakan (\textit{ParticipantObjectTypeCode}, \textit{ParticipantObjectTypeCodeRole}, \textit{ParticipantObjectIDTypeCode}, dan \textit{ParticipantObjectID}) (\cite{iso27789_2021}).

Di luar konten wajib tersebut, terdapat pula sejumlah bidang bersifat opsional (\textit{optional}) maupun kondisional yang dapat ditambahkan sesuai kebutuhan spesifik dari masing-masing jenis event. Sebagai contoh, pada event akses (\textit{access events}) yang mencakup pembuatan, pembacaan, pembaruan, dan penghapusan data, bidang \textit{EventActionCode} dibatasi pada empat nilai tertentu (\textit{Create}, \textit{Read}, \textit{Update}, \textit{Delete}) dan objek yang diaudit secara spesifik merepresentasikan data pasien (\cite{iso27789_2021}). Sebaliknya, pada event kueri (\textit{query events}), dibutuhkan bidang tambahan yang tidak diperlukan pada event akses, yaitu identitas pihak yang mengajukan kueri (\textit{Questioner related}), pihak yang merespons kueri (\textit{Question ahead related}), serta pihak lain yang turut terlibat (\textit{Alternative participant related}), dan bidang \textit{ParticipantObjectQuery} yang pada event kueri berubah sifatnya dari opsional menjadi kondisional mandatory karena perlu menyimpan isi kueri yang sesungguhnya diajukan (\cite{iso27789_2021}). Bidang opsional lain yang lazim ditambahkan pada berbagai jenis event antara lain indikator keberhasilan tindakan (\textit{EventOutcomeIndicator}), peran pengguna saat melakukan tindakan (\textit{RoleIDCode}), tujuan penggunaan data (\textit{PurposeOfUse}), maupun tingkat sensitivitas objek yang diakses (\textit{ParticipantObjectSensitivity}) (\cite{iso27789_2021}).

\subsubsection{Pengumpulan Event}
\label{subsubsec:pengumpulan-event}

Komponen kedua adalah pengumpulan event, yaitu bagaimana setiap peristiwa yang terjadi pada sistem direkam menjadi suatu catatan (\textit{audit record}) dengan format dan isi yang konsisten. Format umum suatu audit record pada rekam medis elektronik sekurang-kurangnya memuat identitas unik dari peristiwa yang terjadi (\textit{EventID}), jenis tindakan yang dilakukan (\textit{EventActionCode}), waktu terjadinya peristiwa (\textit{EventDateTime}), identitas pengguna atau proses yang melakukan tindakan (\textit{UserID}), serta identitas objek data yang menjadi sasaran tindakan tersebut (\textit{ParticipantObjectID}) (\cite{iso27789_2021}).

Pengumpulan event pada praktiknya menghadapi tantangan tersendiri, terutama pada sistem dengan banyak sumber log (\textit{log sources}) yang berbeda. Setiap sumber log cenderung mencatat informasi dengan format dan kelengkapan yang berbeda-beda, sehingga menyulitkan proses menghubungkan (\textit{correlate}) peristiwa yang tercatat oleh satu sumber dengan peristiwa yang tercatat oleh sumber lainnya (\cite{nist80092}).

\subsubsection{Penyimpanan Event}
\label{subsubsec:penyimpanan-event}

Komponen ketiga adalah penyimpanan event, yaitu bagaimana catatan peristiwa yang telah dikumpulkan disimpan dan dikelola secara aman sepanjang siklus hidupnya. Sistem audit trail perlu menyediakan langkah pengamanan yang cukup untuk melindungi audit log dari perusakan (\textit{tampering}), termasuk mengamankan akses terhadap audit record, mengamankan akses terhadap perangkat audit itu sendiri, serta mencatat seluruh tindakan yang dilakukan terhadap audit trail dengan menyertakan waktu, jenis tindakan, dan pelakunya (\cite{iso27789_2021}). Selain itu, audit record juga perlu tunduk pada kebijakan retensi (\textit{retention policy}) yang idealnya disesuaikan dengan masa berlaku data rekam medis yang bersangkutan, mengingat audit record dapat dibutuhkan sebagai bukti dalam jangka waktu yang panjang (\cite{iso27789_2021}).

Pada level infrastruktur, pertimbangan penyimpanan log turut mencakup keseimbangan antara volume data log yang dihasilkan dengan kapasitas penyimpanan yang tersedia, baik untuk penyimpanan daring (\textit{online}) maupun luring (\textit{offline}), serta kebutuhan keamanan atas data tersebut (\cite{nist80092}).

\subsubsection{Pembacaan Event}
\label{subsubsec:pembacaan-event}

Komponen keempat adalah pembacaan event, yaitu bagaimana catatan peristiwa yang telah tersimpan dapat diakses kembali dan dianalisis oleh pihak yang berwenang. Akses terhadap audit data harus dikendalikan secara ketat dan itu sendiri menjadi objek audit (\textit{itself subject to audit}), dengan akses yang idealnya dilakukan melalui sistem informasi yang dapat menegakkan kontrol tersebut, bukan akses langsung terhadap audit trail (\cite{iso27789_2021}). Fasilitas pembacaan audit trail idealnya juga mendukung analisis berdasarkan kombinasi berbagai bidang data yang tercatat, seperti seluruh akses oleh pengguna tertentu, seluruh tindakan penghapusan yang dilakukan oleh pengguna dengan peran tertentu, atau seluruh peristiwa yang melibatkan data pasien tertentu dalam rentang waktu tertentu (\cite{iso27789_2021}).

Kemampuan pembacaan yang fleksibel semacam ini menjadi penting bagi pihak yang bertanggung jawab atas suatu sistem untuk dapat melakukan investigasi maupun evaluasi kepatuhan kebijakan secara efektif, karena analisis log pada dasarnya merupakan proses berkelanjutan yang perlu dilakukan secara berkala, bukan hanya ketika terjadi insiden (\cite{nist80092}).

\subsubsection{Hash-Chaining}

\textit{Hash-chaining} adalah mekanisme di mana setiap catatan dalam
sebuah urutan menyertakan nilai \textit{hash} kriptografis dari catatan
sebelumnya sebagai bagian dari isinya sendiri (\cite{islam2025logstamping}).
Dengan demikian, setiap catatan secara kriptografis terikat pada
seluruh catatan sebelumnya dalam rantai.

Properti utama yang dihasilkan oleh \textit{hash-chaining} adalah
kemampuan deteksi perusakan: apabila isi satu catatan dimodifikasi,
nilai \textit{hash} catatan tersebut akan berubah, yang pada gilirannya
menyebabkan nilai \textit{hash} yang disimpan di catatan sesudahnya
menjadi tidak konsisten (\cite{islam2025logstamping}). Efek domino ini
berlanjut ke seluruh catatan sesudahnya dalam rantai, sehingga
modifikasi pada posisi mana pun dalam rantai akan terdeteksi melalui
pemeriksaan konsistensi nilai \textit{hash}. Mekanisme ini menjadikan
\textit{hash-chaining} sebagai pendekatan yang efektif untuk menegakkan
sifat \textit{append-only} pada sistem pencatatan.

\subsubsection{Tanda Tangan Digital untuk Keaslian dan Non-Repudiation}

Tanda tangan digital merupakan mekanisme kriptografis yang memungkinkan
suatu entitas membuktikan kepemilikan atas sebuah pesan atau dokumen.
Prinsip kerjanya adalah: pengirim menandatangani data menggunakan
\textit{private key} miliknya, menghasilkan tanda tangan yang dapat
diverifikasi oleh siapa pun yang memiliki \textit{public key}
bersesuaian (\cite{shostack2014}).

Dalam konteks audit trail, tanda tangan digital memenuhi dua kebutuhan
sekaligus. \textit{Pertama}, keaslian (\textit{authenticity}): karena
hanya pemegang \textit{private key} yang dapat menghasilkan tanda
tangan yang valid untuk kunci publik tertentu, verifikasi tanda tangan
membuktikan bahwa audit event benar-benar dibangkitkan oleh komponen
yang mengklaim sebagai sumbernya. \textit{Kedua}, non-repudiation:
apabila suatu entitas telah menandatangani sebuah event, entitas
tersebut tidak dapat menyangkal pernah membangkitkan event tersebut
selama \textit{private key}-nya tidak dikompromikan
(\cite{shostack2014}).

Salah satu skema tanda tangan digital yang banyak digunakan untuk
keperluan ini adalah Ed25519, yang berbasis kurva eliptik Edwards25519.
Skema ini bersifat deterministik (tanda tangan yang dihasilkan selalu
sama untuk input yang sama), memiliki ukuran kunci dan tanda tangan
yang kecil (32 byte untuk kunci publik, 64 byte untuk tanda tangan),
dan efisien secara komputasi, menjadikannya pilihan yang sesuai untuk
lingkungan yang membutuhkan verifikasi tanda tangan dalam volume tinggi.

\subsubsection{Rotasi dan Pengarsipan Log}

Rotasi log adalah praktik manajemen log di mana berkas log aktif secara
berkala diganti dengan berkas baru, sementara berkas lama diarsipkan
untuk keperluan retensi (\cite{nist80092}). Praktik ini memiliki tiga
manfaat utama dalam konteks audit trail.

\textit{Pertama}, manajemen ukuran: berkas log yang terus tumbuh tanpa
batas akan menyulitkan pengelolaan dan dapat menghabiskan kapasitas
penyimpanan. Dengan rotasi berkala, setiap berkas arsip memiliki ukuran
yang dapat diprediksi. \textit{Kedua}, efisiensi analisis: investigasi
seringkali berfokus pada rentang waktu tertentu; berkas arsip yang
terorganisasi berdasarkan periode waktu memudahkan pencarian data yang
relevan tanpa harus memindai seluruh riwayat log. \textit{Ketiga},
keamanan penyimpanan: berkas arsip yang telah dirotasi dapat segera
dipindahkan ke media penyimpanan yang lebih aman dan tahan lama,
terpisah dari lingkungan operasional yang lebih rentan terhadap
gangguan (\cite{nist80092}).

Kombinasi antara \textit{hash-chaining} pada level \textit{record},
tanda tangan digital untuk keaslian \textit{event}, dan rotasi log
untuk pengarsipan berkala membentuk tiga lapisan mekanisme yang saling
melengkapi dalam menjamin integritas, keaslian, dan ketersediaan jangka
panjang dari sebuah sistem audit trail.


\subsection{Kebutuhan Sistem Audit Trail}
\label{subsec:kebutuhan-audit-trail}

Berdasarkan definisi, karakteristik, komponen, dan model layanan yang
telah diuraikan pada subbab sebelumnya, dapat diidentifikasi
kebutuhan-kebutuhan yang harus dipenuhi oleh suatu sistem audit trail
agar dapat berfungsi secara efektif sebagai bukti digital dalam proses
audit maupun investigasi insiden keamanan. Kebutuhan tersebut dapat
dikelompokkan menjadi kebutuhan fungsional dan kebutuhan keamanan.

\subsubsection{Kebutuhan Fungsional}

Kebutuhan fungsional berkaitan dengan kemampuan yang harus dimiliki
sistem audit trail untuk menjalankan fungsinya:

\begin{itemize}
    \item \textbf{Pencatatan event yang komprehensif.} Sistem harus
    mampu mencatat seluruh aktivitas yang relevan dari berbagai
    komponen sistem, tidak terbatas pada satu jenis aktivitas saja.
    Cakupan minimal mencakup peristiwa akses dan peristiwa kueri
    terhadap data sensitif (\cite{iso27789_2021}).

    \item \textbf{Format record yang konsisten.} Setiap audit record
    harus memiliki struktur dan atribut yang seragam, memuat
    setidaknya informasi mengenai jenis event, waktu kejadian, sumber
    event, hasil event, serta identitas aktor dan objek yang terlibat
    (\cite{iso27789_2021}).

    \item \textbf{Penyimpanan terpusat dan terdedikasi.} Seluruh
    audit record harus tersimpan pada repositori khusus yang terpisah
    dari data operasional sistem, sehingga dapat diakses secara
    independen untuk keperluan audit (\cite{iso27789_2021}).

    \item \textbf{Kemampuan pencarian dan analisis.} Sistem harus
    menyediakan fasilitas untuk mengambil dan menganalisis audit record
    berdasarkan berbagai kombinasi atribut seperti rentang waktu, jenis
    event, aktor, maupun objek yang terlibat (\cite{iso27789_2021}).

    \item \textbf{Retensi jangka panjang.} Audit record harus
    dipertahankan selama periode yang sesuai dengan kebutuhan organisasi
    dan regulasi yang berlaku, mempertimbangkan bahwa investigasi dapat
    dilakukan jauh setelah kejadian (\cite{iso27789_2021}).
\end{itemize}

\subsubsection{Kebutuhan Keamanan}

Kebutuhan keamanan berkaitan dengan sifat-sifat yang harus dimiliki
audit record agar dapat dipercaya sebagai bukti digital:

\begin{itemize}
    \item \textbf{Integritas.} Audit record yang telah tersimpan harus
    terlindungi dari modifikasi yang tidak sah. Setiap perubahan pada
    audit record harus dapat terdeteksi (\cite{iso27789_2021}).

    \item \textbf{Keaslian (\textit{authenticity}).} Harus tersedia
    mekanisme untuk memverifikasi bahwa audit record benar-benar
    dibangkitkan oleh komponen sistem yang diklaim sebagai sumbernya,
    bukan oleh pihak lain yang berpura-pura menjadi sumber tersebut
    (\cite{iso27789_2021}).

    \item \textbf{Non-repudiation.} Audit record harus dapat digunakan
    untuk membuktikan bahwa suatu aktivitas benar-benar terjadi dan
    dilakukan oleh entitas tertentu, sehingga entitas tersebut tidak
    dapat menyangkal keterlibatannya (\cite{shostack2014},
    \cite{iso27789_2021}).

    \item \textbf{Sifat \textit{append-only}.} Mekanisme penyimpanan
    audit record harus bersifat tambah-saja, artinya entri yang telah
    tercatat tidak dapat diubah maupun dihapus. Perubahan kondisi
    akibat aktivitas berikutnya dicatat sebagai entri baru yang
    terpisah, bukan sebagai pembaruan entri lama (\cite{iso27789_2021}).

    \item \textbf{Akses terkendali.} Akses terhadap audit record harus
    dikendalikan secara ketat dan akses itu sendiri harus menjadi
    objek audit, sehingga tidak ada pihak yang dapat mengakses
    maupun memanipulasi audit record tanpa meninggalkan jejak
    (\cite{iso27789_2021}).
\end{itemize}

Kebutuhan-kebutuhan tersebut bersifat saling melengkapi: sistem yang
hanya memenuhi kebutuhan fungsional tanpa memenuhi kebutuhan keamanan
tidak dapat diandalkan sebagai bukti digital, sedangkan sistem yang
memenuhi kebutuhan keamanan namun tidak memenuhi kebutuhan fungsional
tidak dapat digunakan untuk investigasi yang komprehensif. Kebutuhan
inilah yang akan digunakan sebagai tolok ukur dalam mengevaluasi
kondisi audit trail pada sistem DecMed di Bab III.

\subsection{Mekanisme Integritas dan Non-Repudiation pada Audit Trail}
\label{subsec:mekanisme-integritas}

Untuk memenuhi kebutuhan keamanan yang telah diidentifikasi pada
Subbab~\ref{subsec:kebutuhan-audit-trail}, khususnya integritas,
keaslian, dan non-repudiation, terdapat beberapa mekanisme kriptografis
yang lazim digunakan dalam perancangan sistem audit trail.

% ======================================================================
% CATATAN REFERENSI UNTUK SUBBAB II.2
% ======================================================================
% Key BARU yang perlu kamu tambahkan ke .bib:
%
% \begin{verbatim}
% @misc{iso27789_2021,
%   title  = {ISO 27789:2021 -- Health Informatics -- Audit Trails for
%             Electronic Health Records},
%   author = {{International Organization for Standardization}},
%   year   = {2021},
%   edition = {Second edition},
%   url    = {https://www.iso.org/standard/76028.html}
% }
%
% @techreport{nist80092,
%   author      = {Kent, Karen and Souppaya, Murugiah},
%   title       = {Guide to Computer Security Log Management},
%   institution = {National Institute of Standards and Technology},
%   year        = {2006},
%   number      = {NIST Special Publication 800-92},
%   doi         = {10.6028/NIST.SP.800-92}
% }
% \end{verbatim}
%
% Key yang SUDAH ADA di .bib kamu dan dipakai lagi di sini:
%   - cert2022insiderthreat (sudah diusulkan sebelumnya di bab1)
%   - laksono2021stride
%
% CATATAN PENTING:
%   - Saya sudah baca langsung isi ISO 27789 dan NIST SP 800-92 dari
%     file yang kamu upload (bukan cuma metadata/abstrak), jadi klaim
%     di atas relatif lebih terjamin akurat dibanding sumber-sumber
%     sebelumnya yang cuma saya cari lewat web. NAMUN tetap disarankan
%     kamu baca ulang klausul yang saya rujuk (terutama Clause 6, 7.1,
%     9.2-9.5 pada ISO 27789, dan Section 2.3.1 pada NIST SP 800-92)
%     untuk memastikan parafrase saya akurat.
%   - Struktur "empat komponen audit trail" (penentuan, pengumpulan,
%     penyimpanan, pembacaan) adalah SINTESIS saya sendiri dari kedua
%     standar tersebut, bukan istilah baku yang persis sama di kedua
%     dokumen aslinya -- ISO 27789 memakai istilah "trigger events",
%     "audit record details", dan "secure management of audit data /
%     access to audit data", sedangkan NIST SP 800-92 memakai istilah
%     "generate, transmit, store, analyze". Kalau dosbing lebih suka
%     istilah yang persis mengikuti salah satu standar, beri tahu saya.
%   - Referensi ke \ref{sec:stride} dan \ref{sec:dlt} yang tadinya ada
%     di draf sebelumnya sudah saya hapus bersama kalimat analisisnya,
%     karena bagian ini sekarang murni menyatakan fakta dari literatur
%     saja -- keterkaitan dengan STRIDE/DecMed/DLT akan dibahas di Bab
%     III sesuai arahanmu.

\section{Keamanan Informasi dan Akuntabilitas}
\label{sec:keamanan-informasi}

\subsection{Prinsip Dasar Keamanan Informasi}
\label{subsec:prinsip-keamanan-informasi}

Keamanan informasi secara klasik didefinisikan melalui tiga properti utama yang dikenal sebagai \textit{CIA Triad}: \textit{Confidentiality} (kerahasiaan), \textit{Integrity} (integritas), dan \textit{Availability} (ketersediaan) (\cite{nist80053r5}). Ketiga properti ini membentuk landasan konseptual bagi sebagian besar standar dan kerangka kerja keamanan sistem informasi, termasuk NIST SP 800-53 dan ISO/IEC 27002.

\textit{Confidentiality} merujuk pada jaminan bahwa informasi hanya dapat diakses oleh pihak-pihak yang berwenang, sehingga melindungi data dari pengungkapan yang tidak sah (\cite{nist80053r5}). Dalam konteks RME, kerahasiaan berarti bahwa data klinis dan administratif pasien tidak boleh dapat dibaca oleh pihak mana pun tanpa otorisasi yang sah dari pasien maupun prosedur yang berlaku.

\textit{Integrity} merujuk pada jaminan bahwa informasi tidak dimodifikasi atau dihancurkan dengan cara yang tidak sah, termasuk jaminan ketidakpenyangkalan dan keaslian informasi (\cite{nist80053r5}). Integritas tidak hanya mencakup perlindungan terhadap perubahan isi data, tetapi juga mencakup jaminan bahwa data yang tersimpan mencerminkan kondisi yang sesungguhnya terjadi pada sistem, bukan hasil manipulasi pihak tidak berwenang.

\textit{Availability} merujuk pada jaminan bahwa sistem informasi dan data yang dikelolanya dapat diakses dan digunakan oleh pihak yang berwenang pada saat dibutuhkan (\cite{nist80053r5}). Ketersediaan menjadi relevan tidak hanya bagi fungsionalitas layanan, tetapi juga bagi sistem audit trail: rekaman yang tidak dapat diakses saat dibutuhkan untuk investigasi tidak dapat berfungsi sebagai bukti yang andal.

Ketiga properti tersebut saling berkaitan dan dalam praktiknya sering kali memiliki tegangan satu sama lain. Upaya meningkatkan kerahasiaan melalui enkripsi data, misalnya, dapat mempersulit verifikasi integritas secara langsung. Sebaliknya, mekanisme yang menjamin ketersediaan data secara luas dapat meningkatkan risiko terhadap kerahasiaannya. Ketegangan ini menjadi salah satu alasan mengapa perancangan sistem keamanan perlu dilakukan secara holistik dengan mempertimbangkan ketiga properti secara bersamaan (\cite{shostack2014}).

\subsection{Properti Keamanan yang Diperluas}
\label{subsec:properti-keamanan-diperluas}

Seiring perkembangan sistem informasi yang semakin kompleks, komunitas keamanan informasi mengakui bahwa \textit{CIA Triad} tidak sepenuhnya mencakup seluruh kebutuhan keamanan yang relevan. Beberapa properti tambahan telah diidentifikasi sebagai pelengkap yang diperlukan, khususnya dalam konteks sistem yang melibatkan banyak pihak dengan tingkat kepercayaan yang berbeda (\cite{iso27789_2021}).

\textbf{\textit{Authenticity}} (keaslian) adalah properti yang menjamin bahwa suatu entitas --- baik pengguna, proses, maupun data --- benar-benar merupakan entitas yang diklaim, bukan pihak lain yang berpura-pura (\cite{nist80053r5}). Keaslian mencakup verifikasi identitas aktor yang berinteraksi dengan sistem, maupun verifikasi bahwa data yang diterima berasal dari sumber yang sah dan tidak mengalami perubahan selama pengiriman. Dalam konteks audit trail, keaslian berarti bahwa setiap catatan peristiwa harus dapat dibuktikan berasal dari komponen sistem yang diklaim sebagai sumbernya.

\textbf{\textit{Non-Repudiation}} (ketidakpenyangkalan) adalah properti yang menjamin bahwa suatu entitas tidak dapat menyangkal telah melakukan suatu tindakan tertentu setelah tindakan tersebut terjadi (\cite{nist80053r5}). Non-repudiation merupakan properti yang beroperasi pada lapisan bisnis dan hukum, bukan semata pada lapisan teknis: tujuannya adalah menyediakan bukti yang secara meyakinkan dapat membuktikan keterlibatan suatu entitas dalam suatu peristiwa, sehingga penyangkalan atas peristiwa tersebut tidak dapat diterima. Non-repudiation umumnya diwujudkan melalui kombinasi tanda tangan digital dan sistem pencatatan yang tidak dapat dimanipulasi (\cite{shostack2014}).

\textbf{\textit{Accountability}} (akuntabilitas) adalah properti yang menjamin bahwa tindakan suatu entitas pada suatu sistem dapat ditelusuri secara unik kepada entitas tersebut (\cite{nist80053r5}). Akuntabilitas merupakan properti yang lebih luas dari non-repudiation: non-repudiation berfokus pada pembuktian bahwa suatu tindakan tertentu pernah dilakukan oleh entitas tertentu, sedangkan akuntabilitas mencakup kemampuan sistem untuk menghubungkan setiap tindakan yang terjadi dengan pelakunya secara menyeluruh, tidak hanya pada kasus yang disengketakan. Standar NIST SP 800-53 menetapkan akuntabilitas sebagai prasyarat bagi penerapan tanggung jawab individual (\textit{individual accountability}) pada sistem informasi, yang mensyaratkan bahwa setiap pengguna sistem dapat dimintai pertanggungjawaban atas tindakannya (\cite{nist80053r5}).

\textbf{\textit{Reliability}} (keandalan) dan \textbf{\textit{Trustworthiness}} (keterpercayaan) merupakan properti pendukung yang menjamin bahwa sistem berperilaku secara konsisten sesuai dengan yang dimaksudkan, dan bahwa keluaran yang dihasilkan oleh sistem dapat dipercaya sebagai representasi yang akurat dari kondisi yang sesungguhnya terjadi (\cite{iso27789_2021}).

\subsection{Keterkaitan antara Properti Keamanan dan Audit Trail}
\label{subsec:keterkaitan-properti-audit-trail}

Sistem audit trail tidak hanya merupakan mekanisme untuk memenuhi kebutuhan keamanan tertentu, tetapi juga merupakan komponen yang secara langsung merealisasikan beberapa properti keamanan sekaligus. Keterkaitan antara properti-properti keamanan yang telah diuraikan dan mekanisme audit trail dapat dijabarkan sebagai berikut.

Pertama, audit trail berkontribusi secara langsung terhadap \textit{integrity}. Ketika suatu sistem memiliki rekaman kronologis atas setiap tindakan yang dilakukan terhadap data, rekaman tersebut dapat digunakan untuk memverifikasi bahwa status data saat ini konsisten dengan seluruh tindakan yang pernah tercatat. Apabila terdapat perbedaan antara status data yang teramati dan riwayat tindakan yang tercatat, hal tersebut mengindikasikan kemungkinan adanya modifikasi yang tidak tercatat, yaitu pelanggaran integritas (\cite{iso27789_2021}).

Kedua, audit trail merupakan mekanisme utama untuk mewujudkan \textit{accountability}. Tanpa pencatatan yang menyeluruh atas tindakan yang dilakukan oleh setiap pengguna dan komponen sistem, akuntabilitas hanya bersifat normatif (kebijakan yang menyatakan siapa bertanggung jawab atas apa) tanpa dapat diwujudkan secara teknis (kemampuan sistem untuk membuktikan siapa melakukan apa dan kapan). NIST SP 800-53 secara eksplisit menempatkan keluarga kontrol \textit{Audit and Accountability} (AU) sebagai mekanisme teknis utama untuk merealisasikan akuntabilitas pada sistem informasi (\cite{nist80053r5}).

Ketiga, audit trail yang dilengkapi dengan mekanisme tanda tangan digital merealisasikan properti \textit{non-repudiation}. Apabila setiap catatan peristiwa ditandatangani secara kriptografis oleh komponen yang membangkitkannya, maka komponen tersebut tidak dapat menyangkal pernah membangkitkan peristiwa tersebut selama kunci privatnya tidak dikompromikan (\cite{shostack2014}). Dalam konteks sistem multipartai seperti DecMed, non-repudiation melalui audit trail menjadi penting karena banyak tindakan yang memiliki konsekuensi hukum --- seperti pemberian akses ke data pasien, pembuatan rekam medis, atau pencabutan akses --- perlu dapat dibuktikan secara meyakinkan bahwa tindakan tersebut benar-benar dilakukan oleh pihak yang mengklaim atau yang diklaim melakukannya.

Keempat, audit trail yang bersifat \textit{tamper-evident} dan disimpan secara terpisah dari data operasional berkontribusi terhadap \textit{availability} dalam arti yang lebih luas: ketersediaan bukti. Dalam investigasi insiden keamanan, ketersediaan bukti yang otentik, lengkap, dan tidak dapat dimanipulasi menentukan apakah investigasi tersebut dapat menghasilkan kesimpulan yang meyakinkan atau tidak (\cite{iso27789_2021}).

Hubungan antara properti keamanan dan mekanisme audit trail disarikan pada Tabel~\ref{tab:properti-audit-trail}.

\begin{table}[h]
\centering
\caption{Keterkaitan Properti Keamanan dengan Mekanisme Audit Trail}
\label{tab:properti-audit-trail}
\begin{tabular}{|p{3cm}|p{4cm}|p{5.5cm}|}
\hline
\textbf{Properti Keamanan} & \textbf{Mekanisme Audit Trail} & \textbf{Kontribusi} \\
\hline
\textit{Integrity} &
\textit{Hash-chaining}, penyimpanan \textit{append-only} &
Memungkinkan deteksi modifikasi data yang tidak tercatat melalui inkonsistensi antara status data dan riwayat tindakan. \\
\hline
\textit{Availability} &
Retensi jangka panjang, penyimpanan terdesentralisasi &
Menjamin ketersediaan bukti saat dibutuhkan untuk investigasi, tidak hanya ketersediaan layanan operasional. \\
\hline
\textit{Authenticity} &
Tanda tangan digital pada \textit{audit event} &
Membuktikan bahwa catatan peristiwa berasal dari komponen yang diklaim sebagai sumbernya. \\
\hline
\textit{Non-Repudiation} &
Tanda tangan digital, pencatatan \textit{immutable} &
Mencegah penyangkalan atas tindakan yang telah dilakukan dan tercatat dalam sistem. \\
\hline
\textit{Accountability} &
Pencatatan menyeluruh atas seluruh tindakan dan aktornya &
Menghubungkan setiap tindakan pada sistem dengan identitas pelakunya secara dapat dibuktikan. \\
\hline
\end{tabular}
\end{table}
% 

\section{Kerja DecMed}
\label{sec:decmed}

\subsection{Gambaran Umum dan Latar Belakang DecMed}
\label{subsec:decmed-gambaran-umum}

DecMed adalah kerangka kerja manajemen RME terdesentralisasi yang dikembangkan untuk menjawab kebutuhan pengelolaan data rekam medis yang aman, akuntabel, dan berpusat pada kendali pasien (\cite{DecMed}). Kerangka kerja ini mengintegrasikan tiga teknologi utama: DLT IOTA sebagai lapisan penyimpanan \textit{on-chain} yang \textit{immutable} untuk data kontrol akses, IPFS sebagai penyimpanan \textit{off-chain} untuk berkas klinis terenkripsi berukuran besar, serta PRE sebagai mekanisme kriptografi yang memungkinkan distribusi akses data kepada pihak yang berwenang tanpa mengungkap isi data kepada perantara (\cite{DecMed}). Kombinasi ketiga teknologi ini dirancang agar setiap akses terhadap data RME hanya dapat dilakukan atas persetujuan aktif pasien, dengan durasi dan cakupan akses yang dapat dibatasi oleh pasien sendiri melalui mekanisme CapBAC (\cite{DecMed}).

Arsitektur DecMed mengikuti pola tiga lapis yang telah diuraikan pada Subbab~\ref{subsec:pola-layanan-decmed}. Tiga aplikasi \textit{client} utama, yaitu \textit{Patient Client} (berbasis mobile), \textit{Hospital Client} (berbasis desktop), dan \textit{Ministry Client} (berbasis desktop), berfungsi sebagai antarmuka bagi masing-masing peran aktor dalam sistem. Seluruh aplikasi \textit{client} dibangun menggunakan Tauri dengan \textit{frontend} SvelteKit/TypeScript, sementara logika \textit{backend} ditulis dalam Rust (\cite{DecMed}). Identitas kriptografi pengguna dikelola secara \textit{non-custodial} menggunakan skema BIP39, di mana setiap pengguna menyimpan kunci privat mereka sendiri dalam bentuk \textit{mnemonic} 12 kata tanpa bergantung pada pihak ketiga (\cite{DecMed}).

\subsection{Mekanisme Kontrol Akses pada DecMed}
\label{subsec:decmed-kontrol-akses}

Mekanisme kontrol akses pada DecMed dibangun di atas model CapBAC yang memanfaatkan \textit{smart contract} Move pada jaringan IOTA sebagai repositori \textit{capability} yang \textit{immutable} (\cite{DecMed}). Ketika seorang pasien memberikan akses kepada tenaga medis atau tenaga administratif, pasien membuat entri \textit{capability} pada \textit{smart contract} yang memuat atribut seperti identitas penerima akses, jenis akses yang diberikan (\textit{read} atau \textit{update}), dan waktu kedaluwarsa akses. Entri ini tersimpan secara \textit{on-chain} dan tidak dapat diubah oleh pihak lain selain pasien selaku pemilik data (\cite{DecMed}).

Setiap kali tenaga medis mengajukan permintaan akses terhadap data RME, PRE Server memverifikasi keberadaan entri \textit{capability} yang bersesuaian pada \textit{smart contract} sebelum melakukan \textit{re-encryption} data. Proses ini memastikan bahwa akses hanya diberikan kepada pihak yang memiliki \textit{capability} aktif dan belum kedaluwarsa (\cite{DecMed}). Selain memberikan akses, pasien juga dapat mencabut akses (\textit{revoke access}) sewaktu-waktu, yang menyebabkan entri \textit{capability} pada \textit{smart contract} diperbarui dengan status \textit{revoked} (\cite{DecMed}).

\subsection{Keterbatasan Audit Trail pada DecMed}
\label{subsec:decmed-keterbatasan-audit}

DecMed mengklaim menghasilkan mekanisme kontrol akses yang didukung oleh audit trail yang tahan terhadap perusakan (\textit{tamper-evident audit trails}) (\cite{DecMed}). Klaim tersebut didasarkan pada mekanisme bahwa setiap akses yang diberikan oleh pasien dicatat sebagai entri \textit{capability} pada penyimpanan \textit{on-chain} di jaringan IOTA, dan sifat \textit{immutability} dari penyimpanan \textit{on-chain} tersebut menjamin bahwa entri tidak dapat diubah oleh pihak lain, dengan setiap permintaan akses divalidasi terhadap keberadaan entri yang bersesuaian melalui \textit{smart contract} (\cite{DecMed}). Peristiwa akses, pemberian persetujuan (\textit{consent}), dan pelanggaran kebijakan pada DecMed dinyatakan tercatat secara \textit{tamper-evident}, sehingga menyediakan audit trail yang dapat diandalkan untuk mendukung akuntabilitas lintas batas organisasi (\cite{DecMed}).

Meskipun demikian, apabila ditinjau terhadap karakteristik audit trail yang baik sebagaimana diuraikan pada Subbab~\ref{subsec:karakteristik-audit-trail}, mekanisme pencatatan yang tersedia pada DecMed memiliki sejumlah keterbatasan mendasar. Pertama, sebagian besar mekanisme yang tercatat pada jaringan IOTA merupakan entri \textit{capability} atau izin akses (\textit{access grant}), yaitu metadata yang menyatakan siapa berhak mengakses data apa dan sampai kapan, bukan catatan atas peristiwa akses yang benar-benar terjadi (\textit{access event}), sehingga sistem belum dapat merekonstruksi riwayat siapa yang benar-benar membaca atau mengubah data RME pada waktu tertentu. Kedua, sejumlah operasi pada DecMed, seperti pencabutan akses (\texttt{revoke\_access}) dan pembaruan data administratif maupun klinis, menyebabkan entri lama pada state tergantikan oleh entri baru alih-alih ditambahkan sebagai riwayat baru yang bersifat \textit{append-only}, yang bertentangan dengan karakteristik dasar audit trail mengenai integritas dan \textit{non-repudiation} (\cite{DecMed}). Keterbatasan-keterbatasan inilah yang menjadi motivasi utama pengembangan sistem audit trail terintegrasi pada Tugas Akhir ini, sebagaimana akan dianalisis lebih lanjut pada Bab III.

\section{Distributed Ledger Technology dan Immutability}
\label{sec:dlt}

\subsection{DLT sebagai Basis Pencatatan Tidak Berubah}
\label{subsec:dlt-umum}

Distributed Ledger Technology (DLT) merupakan sistem terdesentralisasi yang memungkinkan data direkam dan disimpan pada sejumlah \textit{node} dalam suatu jaringan, sehingga seluruh salinan \textit{ledger} tetap tersinkronisasi dan konsisten satu sama lain (\cite{DecMed}). Karakteristik desentralisasi ini menghilangkan kebutuhan akan otoritas terpusat, sehingga meningkatkan transparansi dalam sistem, dan menjadikan DLT banyak digunakan di berbagai sektor, termasuk keuangan, \textit{supply chain}, kesehatan, dan energi (\cite{DecMed}).

Salah satu karakteristik utama DLT yang relevan bagi kebutuhan audit trail adalah \textit{immutability}, yaitu sifat data yang telah tersimpan pada \textit{ledger} tidak dapat diubah maupun dihapus setelah tercatat. Karakteristik ini menjadikan DLT sebagai salah satu teknologi yang cocok untuk diintegrasikan dengan mekanisme kontrol akses berbasis \textit{capability} (CapBAC), karena DLT dapat berfungsi sebagai media penyimpanan bagi \textit{capability} yang membutuhkan jaminan integritas (\cite{DecMed}).

\subsection{IOTA}
\label{subsec:iota}

IOTA merupakan salah satu jenis DLT yang mengadopsi struktur jaringan berbasis \textit{Directed Acyclic Graph} (DAG). Struktur jaringan DAG memungkinkan IOTA mendukung \textit{throughput} yang tinggi, karena mampu memvalidasi transaksi baru secara lebih cepat dibandingkan \textit{blockchain} konvensional seperti Bitcoin maupun Ethereum (\cite{DecMed}). Pada versi awalnya (IOTA 2.0), IOTA tidak membebankan biaya transaksi (\textit{zero-cost transactions}), sehingga cocok digunakan untuk transaksi mikro maupun pada lingkungan ringan seperti perangkat Internet of Things (IoT). Namun, pada IOTA Rebased, model transaksi tanpa biaya tersebut tidak lagi berlaku, dan setiap transaksi dikenakan biaya (\textit{gas fee}) (\cite{DecMed}).

IOTA Rebased menandai evolusi signifikan dari protokol IOTA, bertransisi dari struktur Tangle yang mendukung transaksi tanpa biaya dan bergantung pada \textit{coordinator} terpusat, menuju \textit{object-based ledger} yang dibangun di atas DAG. Berbeda dari IOTA 2.0 yang berbasis model \textit{Unspent Transaction Output} (UTXO), IOTA Rebased mengadopsi DAG berbasis objek yang terinspirasi dari arsitektur Sui Network, sehingga memungkinkan pemrosesan transaksi secara paralel dan mendukung pola transaksi yang kompleks melalui kontrol akses \textit{fine-grained} (\cite{DecMed}). \textit{Ledger} pada IOTA Rebased membedakan antara \textit{shared objects}, yang dapat diakses dan dimodifikasi oleh banyak entitas berdasarkan aturan yang telah ditentukan, dengan \textit{owned objects}, yang sepenuhnya dikontrol oleh satu entitas (\cite{DecMed}). Konsensus pada IOTA Rebased ditenagai oleh protokol Mysticeti, sebuah sistem \textit{Byzantine Fault Tolerant} (BFT) yang mengimplementasikan \textit{Delegated Proof-of-Stake} (DPoS) untuk mencapai \textit{throughput} tinggi dan \textit{latency} rendah, dengan \textit{finality} dalam waktu sekitar 500 milidetik dan mendukung lebih dari 50.000 transaksi per detik (\cite{DecMed}).

Untuk menyederhanakan interaksi pengguna dengan jaringan IOTA, tersedia IOTA Gas Station yang memungkinkan transaksi bersponsor (\textit{sponsored transactions}), sehingga pengguna tidak perlu mengelola \textit{token} untuk membayar biaya transaksi (\cite{DecMed}). Arsitektur IOTA Gas Station terdiri atas tiga komponen, yaitu Gas Station Server sebagai komponen utama yang menyediakan API untuk mereservasi \textit{gas object} dan mengeksekusi transaksi, Gas Station Storage sebagai penyimpanan persisten untuk \textit{gas object} dan status transaksi bersponsor, serta Key Store Manager (KSM) yang mengelola kunci untuk menandatangani transaksi, baik melalui layanan eksternal (\textit{External Key Management Service}) maupun penyimpanan kunci dalam memori (\cite{DecMed}).

% Tambahkan setelah paragraf tentang IOTA Gas Station (paragraf terakhir subbab IOTA)

Setiap entitas yang berinteraksi dengan jaringan IOTA Rebased diidentifikasi
melalui \textit{IOTA address}, yaitu pengenal 32-byte yang secara kriptografis
diturunkan dari kunci publik entitas tersebut (\cite{iotadocs_keys_addresses}).
Pada skema penandatanganan bawaan Ed25519, \textit{address} diperoleh dengan
menghitung \textit{hash} BLAKE2b-256 dari kunci publik Ed25519 secara langsung,
tanpa awalan \textit{flag}; sedangkan untuk skema penandatanganan lainnya seperti
Secp256k1 dan Secp256r1, satu byte \textit{flag} skema penandatanganan
dikombinasikan terlebih dahulu dengan kunci publik sebelum dihitung nilai
\textit{hash}-nya (\cite{iotadocs_keys_addresses}). Hasil \textit{hash} yang
dihasilkan direpresentasikan sebagai \textit{string} heksadesimal dengan awalan
\texttt{0x}, misalnya \texttt{0x02a212de6a9dfa3a69e22387acfbafbb1a9e591bd9d636e7895dcfc8de05f331}
(\cite{iotadocs_keys_addresses}). Setiap \textit{address} bersifat unik dan
secara kriptografis terikat pada kunci privat yang bersesuaian: hanya pemegang
kunci privat tersebut yang dapat mengotorisasi transaksi yang dikirimkan atas
nama \textit{address} bersangkutan, tanpa memerlukan inisialisasi \textit{on-chain}
terlebih dahulu (\cite{iotadocs_exchange_integration}).

Dalam arsitektur DecMed, \textit{IOTA address} berperan ganda: sebagai
identitas \textit{on-chain} yang digunakan untuk kepemilikan objek dan
penandatanganan transaksi, sekaligus sebagai pengenal aktor dalam mekanisme
kontrol akses berbasis \textit{capability} (CapBAC). Setiap pasien, tenaga
fasyankes, dan administrator memiliki \textit{IOTA address} masing-masing yang
diturunkan dari \textit{recovery phrase} BIP39 mereka (\cite{DecMed}). Ketika
suatu \textit{capability} akses dibuat oleh pasien untuk diberikan kepada tenaga
medis tertentu, \textit{address} tenaga medis tersebut dicantumkan secara
eksplisit sebagai penerima \textit{capability}, sehingga validasi otorisasi pada
\textit{smart contract} Move dapat memverifikasi kesesuaian antara pengirim
transaksi dan penerima \textit{capability} yang tercatat (\cite{DecMed}).
Peran ini menjadikan \textit{IOTA address} sebagai basis identifikasi aktor
yang relevan untuk dicatat pada setiap \textit{audit event}, sebagaimana
digunakan pada atribut \texttt{actor\_id} dalam skema audit trail yang
dikembangkan pada Tugas Akhir ini.

\section{Arsitektur Layanan Pendukung dalam Sistem Terdesentralisasi}
\label{sec:arsitektur-layanan-pendukung}

\subsection{Pola Layanan Terpisah dalam Arsitektur DecMed}
\label{subsec:pola-layanan-decmed}

Arsitektur DecMed dirancang mengikuti pola tiga lapis (\textit{three-tier architecture}), yang terdiri atas \textit{Client Layer}, \textit{Proxy Layer}, dan \textit{Storage Layer} (\cite{DecMed}). \textit{Client Layer} berfungsi sebagai antarmuka utama bagi seluruh aktor untuk berinteraksi dengan sistem, dengan tiga aplikasi klien berbeda yang disesuaikan dengan kebutuhan masing-masing peran. \textit{Proxy Layer} berperan sebagai \textit{middleware} keamanan yang menegakkan kebijakan kontrol akses dan mengelola operasi kriptografi seperti Proxy Re-Encryption (PRE), sedangkan \textit{Storage Layer} terbagi menjadi penyimpanan \textit{on-chain} yang menjamin keutuhan \textit{capability} akses, dan penyimpanan \textit{off-chain} yang menangani data klinis terenkripsi dalam volume besar (\cite{DecMed}).

Lapisan \textit{Proxy} dan \textit{Storage} pada DecMed diimplementasikan melalui beberapa layanan (\textit{server}) yang terpisah satu sama lain, masing-masing dengan tanggung jawab yang spesifik. Layanan \textit{PRE Server} bertanggung jawab menangani proses \textit{re-encryption}, di mana setiap kali menerima permintaan data, \textit{server} tersebut terlebih dahulu memverifikasi entri \textit{capability} pada jaringan IOTA sebelum melakukan \textit{re-encryption} terhadap data yang diambil dari IPFS untuk pemohon, dengan Redis digunakan untuk menyimpan data sementara seperti \textit{nonce} guna mencegah \textit{replay attack} (\cite{DecMed}). Layanan \textit{Gas Station Server} bertanggung jawab mensponsori transaksi pada jaringan IOTA, sehingga aktor yang tidak memiliki kepentingan finansial dapat berinteraksi dengan \textit{ledger} tanpa perlu memiliki \textit{token} (\cite{DecMed}). Kedua layanan tersebut berkomunikasi dengan jaringan IOTA maupun IPFS sesuai kebutuhannya masing-masing, sementara lapisan penyimpanan sendiri terbagi menjadi IOTA Tangle sebagai \textit{trust anchor} yang immutable untuk kebijakan kontrol akses (CapBAC), dan IPFS Cluster sebagai penyimpanan \textit{content-addressable} yang terdesentralisasi untuk berkas RME terenkripsi (\cite{DecMed}).

\subsection{Proxy Re-Encryption dan IPFS}
\label{subsec:pre-ipfs}

Proxy Re-Encryption (PRE) memanfaatkan \textit{proxy} yang bersifat \textit{semi-trusted} untuk mentransformasikan suatu \textit{ciphertext} menjadi \textit{ciphertext} baru yang isinya tetap sama, namun hanya dapat didekripsi menggunakan \textit{secret key} milik subjek yang dituju (\cite{DecMed}). \textit{Ciphertext} awal dihasilkan dari enkripsi data menggunakan \textit{public key} milik \textit{delegator} (pemilik data), dan dapat disimpan pada suatu lokasi penyimpanan seperti \textit{cloud storage} maupun sistem DLT. Ketika \textit{delegatee} (penerima akses) meminta akses terhadap data terenkripsi tersebut, \textit{delegator} membuat \textit{re-encryption key} menggunakan \textit{secret key} miliknya dan \textit{public key} milik \textit{delegatee}; \textit{re-encryption key} tersebut dikirimkan kepada \textit{proxy}, yang kemudian menggunakannya bersama \textit{ciphertext} awal dan \textit{public key} kedua belah pihak untuk menghasilkan \textit{ciphertext} baru yang hanya dapat didekripsi oleh \textit{delegatee} (\cite{DecMed}). Meskipun PRE membatasi kemampuan \textit{proxy} untuk mengetahui isi data asli (\textit{plaintext}), arsitektur ini tetap memusatkan fungsi \textit{re-encryption} pada satu entitas \textit{semi-trusted}, sehingga menciptakan \textit{single point of failure}: apabila \textit{proxy} mengalami gangguan atau menjadi target serangan \textit{denial-of-service}, seluruh proses \textit{re-encryption} akan terhenti dan alur otorisasi maupun distribusi kunci menjadi tidak tersedia bagi seluruh subjek yang bergantung padanya (\cite{DecMed}).

InterPlanetary File System (IPFS) adalah suatu sistem berkas terdistribusi \textit{peer-to-peer} yang bertujuan menghubungkan seluruh perangkat komputasi ke dalam satu sistem berkas (\cite{benet2014ipfscontentaddressed}). Berbeda dengan protokol Hypertext Transfer Protocol (HTTP) konvensional yang menggunakan pengalamatan berbasis lokasi (\textit{location-based addressing}), IPFS menggunakan pengalamatan berbasis konten (\textit{content-based addressing}), yaitu data diidentifikasi dan diambil berdasarkan \textit{hash} kriptografisnya yang disebut \textit{Content Identifier} (CID) (\cite{DecMed}). Ketika suatu berkas ditambahkan ke jaringan IPFS, berkas tersebut dipecah menjadi bagian-bagian kecil, di-\textit{hash} secara kriptografis, dan diberikan CID yang unik sebagai penanda permanen dan \textit{immutable} dari data tersebut; apabila isi berkas berubah, \textit{hash} kriptografisnya turut berubah dan menghasilkan CID yang sepenuhnya baru, sehingga karakteristik \textit{immutability} ini menjadikan IPFS sangat kompatibel dengan DLT dalam menjaga integritas data (\cite{DecMed}). Dalam konteks penyimpanan data kesehatan, IPFS memungkinkan penyimpanan \textit{off-chain} untuk berkas medis terenkripsi berukuran besar, sementara hanya CID berukuran kecil yang perlu disimpan pada \textit{ledger} yang \textit{immutable} (\textit{on-chain}), sehingga mengatasi keterbatasan ukuran penyimpanan yang umum ditemukan pada jaringan DLT maupun \textit{blockchain} tanpa mengorbankan arsitektur terdesentralisasi yang bebas dari \textit{single point of failure} (\cite{DecMed}).

% Tambahkan setelah paragraf tentang IPFS (paragraf terakhir tentang IPFS
% yang membahas CID dan penyimpanan off-chain)

Meskipun IPFS menyediakan mekanisme pengalamatan berbasis konten yang menjamin
identifikasi data secara permanen melalui CID, keberadaan suatu berkas pada
jaringan IPFS tidak secara otomatis berarti berkas tersebut akan tersedia
selamanya. Setiap \textit{node} IPFS hanya menyimpan data yang secara eksplisit
dimintanya atau yang pernah diaksesnya, dan \textit{node} tersebut dapat
menjalankan proses \textit{Garbage Collection} secara berkala untuk menghapus
data yang tidak lagi direferensikan secara aktif demi membebaskan kapasitas
penyimpanan (\cite{benet2014ipfscontentaddressed}). Akibatnya, sebuah berkas
yang diunggah ke IPFS tanpa penanganan khusus berpotensi menjadi tidak dapat
diakses apabila seluruh \textit{node} yang menyimpan salinannya memutuskan
untuk menghapusnya melalui \textit{Garbage Collection}.

Untuk mengatasi persoalan ini, IPFS menyediakan mekanisme \textit{pinning},
yaitu penanda eksplisit yang diberikan kepada suatu CID untuk memberitahu
\textit{node} bahwa data yang dirujuk oleh CID tersebut harus dipertahankan
dan tidak boleh dihapus oleh proses \textit{Garbage Collection}
(\cite{benet2014ipfscontentaddressed}). Sebuah berkas yang telah di-\textit{pin}
pada suatu \textit{node} akan tetap tersimpan pada \textit{node} tersebut
selama penanda \textit{pin} tidak dicabut secara eksplisit, terlepas dari
seberapa jarang berkas tersebut diakses. Dalam konteks penyimpanan data yang
membutuhkan retensi jangka panjang, seperti berkas log audit trail, mekanisme
\textit{pinning} menjadi prasyarat agar data yang diunggah ke IPFS dapat
diandalkan ketersediaannya dalam jangka waktu yang panjang. IPFS Cluster,
yang digunakan oleh DecMed untuk penyimpanan data klinis, menyediakan
fungsionalitas \textit{pinning} terdistribusi yang memastikan CID yang
di-\textit{pin} direplikasi dan dipertahankan pada seluruh \textit{node}
dalam klaster, sehingga ketidaktersediaan satu \textit{node} tidak
menyebabkan hilangnya data (\cite{DecMed}).

\section{Threat Modeling dengan STRIDE}
\label{sec:stride}

\subsection{Konsep dan Kategori STRIDE}
\label{subsec:konsep-kategori-stride}

\textit{Threat modeling} (pemodelan ancaman) merupakan proses terstruktur untuk mengidentifikasi, mengenumerasi, dan memprioritaskan ancaman keamanan pada suatu sistem sebelum ancaman tersebut benar-benar dieksploitasi (\cite{shostack2014}). Proses ini umumnya terdiri atas empat pertanyaan pokok: apa yang sedang dibangun, apa yang dapat salah, apa yang perlu dilakukan terhadap hal tersebut, dan apakah pekerjaan tersebut telah dilakukan dengan cukup baik (\cite{shostack2014}).

STRIDE merupakan salah satu kerangka kerja \textit{threat modeling} yang dikembangkan oleh Loren Kohnfelder dan Praerit Garg di Microsoft, dengan nama yang dibentuk dari huruf pertama masing-masing kategori ancaman (\cite{shostack2014}). Kerangka kerja ini dirancang untuk membantu pengembang perangkat lunak mengidentifikasi jenis-jenis serangan yang cenderung dialami suatu sistem. Setiap kategori STRIDE merupakan kebalikan dari properti keamanan yang diinginkan pada suatu sistem, sebagaimana ditunjukkan pada Tabel~\ref{tab:stride-kategori}.

\begin{table}[h]
\centering
\caption{Kategori Ancaman STRIDE dan Properti Keamanan yang Dilanggar}
\label{tab:stride-kategori}
\begin{tabular}{|p{2.5cm}|p{3cm}|p{7cm}|}
\hline
\textbf{Kategori} & \textbf{Properti yang Dilanggar} & \textbf{Definisi} \\
\hline
Spoofing & Authentication & Berpura-pura menjadi sesuatu atau seseorang selain diri sendiri. \\
\hline
Tampering & Integrity & Memodifikasi sesuatu pada disk, jaringan, atau memori. \\
\hline
Repudiation & Non-repudiation & Mengklaim tidak melakukan sesuatu atau tidak bertanggung jawab atas suatu tindakan. \\
\hline
Information Disclosure & Confidentiality & Memberikan informasi kepada pihak yang tidak berwenang untuk melihatnya. \\
\hline
Denial of Service & Availability & Menyerap sumber daya yang dibutuhkan untuk menyediakan layanan. \\
\hline
Elevation of Privilege & Authorization & Memungkinkan seseorang melakukan sesuatu yang tidak diizinkan untuk dilakukan. \\
\hline
\end{tabular}
\end{table}

Kategori \textit{Repudiation} memiliki karakteristik yang berbeda dari kategori lainnya: ancaman ini beroperasi pada lapisan bisnis, bukan semata pada lapisan teknis jaringan atau aplikasi (\cite{shostack2014}). Penting untuk dipahami bahwa repudiasi tidak selalu merupakan tindakan jahat, misalnya seseorang dapat secara jujur menyatakan tidak melakukan sesuatu karena kegagalan teknologi atau proses. Oleh karena itu, pertanyaan kunci bagi perancang sistem adalah: \textit{bukti apa yang dimiliki sistem untuk membantah atau membuktikan suatu klaim?} Apabila sistem tidak memiliki log, tidak menyimpan log, atau tidak dapat menganalisis log, ancaman repudiasi menjadi sangat sulit untuk dibantah (\cite{shostack2014}). Hal ini menegaskan keterkaitan langsung antara kategori \textit{Repudiation} pada STRIDE dengan kebutuhan mekanisme audit trail pada suatu sistem.

\subsection{Dekomposisi Sistem dan Data Flow Diagram dalam Threat Modeling}
\label{subsec:dfd-threat-modeling}

Sebelum mengidentifikasi ancaman menggunakan STRIDE, suatu sistem perlu didekomposisi terlebih dahulu untuk memperoleh pemahaman menyeluruh mengenai cara kerja sistem dan bagaimana sistem tersebut berinteraksi dengan entitas eksternal (\cite{owasp_threat_modeling}). Dekomposisi ini umumnya mencakup identifikasi deskripsi sistem, dependensi eksternal, titik masuk (\textit{entry points}), aset yang perlu dilindungi, dan tingkat kepercayaan (\textit{trust level}) dari setiap entitas yang terlibat.

\textit{Data Flow Diagram} (DFD) merupakan alat utama yang digunakan dalam tahap dekomposisi ini. DFD menggambarkan bagaimana data mengalir di antara elemen-elemen sistem, sehingga memungkinkan analis untuk mengidentifikasi di mana ancaman keamanan paling mungkin muncul (\cite{owasp_threat_modeling}). DFD dalam konteks \textit{threat modeling} terdiri atas empat elemen dasar:

\begin{itemize}
    \item (\textit{Process}): elemen yang menerima, mengolah, dan mengirimkan data, biasanya direpresentasikan sebagai lingkaran atau persegi panjang dengan sudut membulat.
    \item (\textit{Data Store}): tempat data disimpan secara persisten; direpresentasikan sebagai dua garis paralel.
    \item (\textit{Data Flow}): perpindahan data antar elemen; direpresentasikan sebagai panah berlabel.
    \item (\textit{External Entity}): aktor atau sistem di luar kendali sistem yang dianalisis; direpresentasikan sebagai persegi panjang.
\end{itemize}

Elemen tambahan yang krusial dalam \textit{threat modeling} adalah \textbf{batas kepercayaan} (\textit{trust boundary}), yang memisahkan bagian-bagian sistem dengan tingkat kepercayaan berbeda dan direpresentasikan sebagai garis putus-putus. Ancaman paling signifikan umumnya muncul pada aliran data yang melintasi batas kepercayaan ini, karena pada titik itulah asumsi mengenai kepercayaan terhadap data maupun pengirimnya berubah (\cite{owasp_threat_modeling}).



\section{Penelitian Terkait}
\label{sec:penelitian-terkait}

Salah satu contoh sistem audit trail yang secara khusus dibangun dengan memanfaatkan kombinasi DLT dan IPFS adalah LogStamping (\cite{islam2025logstamping}). Sistem tersebut menggunakan Ethereum sebagai platform \textit{blockchain} dengan konsensus \textit{Proof of Authority}, di mana log dikelompokkan menjadi beberapa \textit{chunk} berdasarkan jumlah baris maupun interval waktu, kemudian dihitung \textit{hash} SHA-256 dari setiap \textit{chunk} tersebut dan disimpan pada \textit{blockchain}, sementara berkas log asli yang telah terenkripsi dan terverifikasi diarsipkan secara terpisah pada IPFS untuk menjamin keberadaannya secara \textit{tamper-evident} dan persisten (\cite{islam2025logstamping}). Pendekatan tersebut menghasilkan audit trail yang ringan (\textit{lightweight}) dan dapat diverifikasi, karena \textit{blockchain} hanya menyimpan \textit{hash} sebagai bukti integritas alih-alih menyimpan seluruh data log secara langsung, sedangkan mekanisme verifikasi bekerja dengan menghitung ulang \textit{hash} dari suatu kelompok log dan membandingkannya dengan \textit{hash} yang tersimpan pada \textit{blockchain} untuk mendeteksi adanya perubahan (\cite{islam2025logstamping}). Hasil pengujian pada sistem tersebut menunjukkan tingkat deteksi perusakan (\textit{tamper detection}) sebesar 100\% terhadap simulasi injeksi log palsu, serta pengurangan kebutuhan penyimpanan hingga 92\% dibandingkan penyimpanan log mentah secara langsung pada \textit{blockchain} (\cite{islam2025logstamping}).